\begin{question}
\noindent An astronomer is using a satellite dish that has a circular cross-section when viewed from the side. A support strut is attached at an external point $P$, with two straight cables $PA$ and $PB$ running from $P$ to touch the circular rim at points $A$ and $B$, where $O$ is the centre of the circle.

Both $PA$ and $PB$ are tangents to the circle, and $\angle AOB = 130^\circ$.

\begin{center}
\begin{tikzpicture}[scale=1.8]
    % Define the center of the circle
    \coordinate (O) at (0,0);

    % Draw the circle
    \draw (O) circle (2cm);
    \fill[black] (O) circle (1.2pt);
    \node at (0,-0.25) {$O$};

    % Place points A and B on the circumference symmetric about the vertical axis
    % angle AOB = 130 degrees, so each is 65 degrees from the downward vertical axis
    \coordinate (A) at ($(O) + (155:2cm)$);
    \coordinate (B) at ($(O) + (25:2cm)$);

    % Draw tangent lines from A and B meeting at P
    \path[name path=tanA] (A) -- ($(A)!4!90:(O)$);
    \path[name path=tanB] (B) -- ($(B)!4!-90:(O)$);
    \path [name intersections={of=tanA and tanB, by={P}}];

    \draw (A) -- (P) -- (B);
    \draw (O) -- (A);
    \draw (O) -- (B);

    % Right angle marks at A and B
    \draw ($(A)!0.18cm!(O)$) -- ($($(A)!0.18cm!(O)$)!0.18cm!90:(O)$) -- ($(A)!0.18cm!90:(O)$);
    \draw ($(B)!0.18cm!(O)$) -- ($($(B)!0.18cm!(O)$)!0.18cm!-90:(O)$) -- ($(B)!0.18cm!-90:(O)$);

    % Labels
    \node at (A) [above left] {$A$};
    \node at (B) [above right] {$B$};
    \node at (P) [below] {$P$};

    % Mark angle AOB
    \pic [draw, "$130^\circ$", angle eccentricity=1.4, angle radius = 0.5cm] {angle = B--O--A};

    % Mark angle APB
    \pic [draw, "$x$", angle eccentricity=1.5, angle radius = 0.7cm] {angle = A--P--B};
\end{tikzpicture}
\end{center}

\noindent Work out the size of angle $APB$, marked $x$ on the diagram.

\vspace{4cm}

\begin{flushright}
\answerequnit{x}{$^\circ$}{3}
\end{flushright}
\end{question}
